\section{Conclusión}

El presente hito consolidó la capa de backend del sistema como una API REST completa, 
estructurada y coherente. Los seis módulos funcionales fueron implementados siguiendo una 
arquitectura uniforme de repositorio, DTO y controlador, lo que garantiza 
predictibilidad en el comportamiento de cada servicio y facilita su consumo desde el 
frontend en los hitos siguientes.

Un aspecto relevante que este hito pone de manifiesto es que el desarrollo del backend 
no es un proceso que ocurre una sola vez y queda cerrado. Las migraciones generadas 
con Alembic durante este hito evidencian que el esquema de base de datos evolucionó 
de forma iterativa: ajustes en campos, eliminación de columnas innecesarias, 
correcciones en relaciones de cascada y la formalización de tablas intermedias como 
modelos ORM explícitos.

Desde el punto de vista del diseño, los módulos más complejos, son correspondidos a Dispositivos y 
Préstamos, los cuales ilustran cómo la lógica de negocio puede encapsularse limpiamente en el 
repositorio, dejando los controladores livianos y enfocados en la orquestación.
El módulo de Préstamos, en particular, demuestra cómo una operación aparentemente 
simple como registrar una devolución involucra actualizaciones coordinadas en múltiples 
tablas dentro de una misma transacción, haciendo evidente la importancia de un diseño 
de datos sólido desde el inicio del proyecto.

Con este backend en funcionamiento, el sistema cuenta con todos los endpoints necesarios 
para que los hitos de frontend puedan avanzar de forma modular e independiente, 
consumiendo una API estable, tipada y documentada.